% Bagian: incompleteness
% Bab: theories-computability
% Subbagian: complete-decidable

\documentclass[../../../include/open-logic-section]{subfiles}

\begin{document}

\olfileid[id]{inc}{tcp}{cdc}

\olsection{Teori Lengkap yang \printtoken{S}{axiomatizable} Dapat Diputuskan}

Suatu teori disebut {\em lengkap} jika, untuk setiap kalimat~$!A$, $!A$
atau $\lnot !A$ dapat dibuktikan.

\begin{lem}
Andaikan suatu teori~$\Th{T}$ lengkap dan !!{axiomatizable}. Maka $\Th{T}$
dapat diputuskan.
\end{lem}

\begin{proof}
Andaikan $\Th{T}$ lengkap dan $A$ merupakan himpunan aksioma yang dapat
dihitung. Jika $\Th{T}$ tak konsisten, teori itu jelas dapat dihitung.
(Algoritmenya: ``cukup jawab ya.'') Jadi, kita dapat mengasumsikan bahwa
$\Th{T}$ juga konsisten.

Untuk memutuskan apakah suatu kalimat~$!A$ berada dalam~$\Th{T}$ atau tidak,
carilah secara bersamaan !!a{derivation} atas~$!A$ dari~$\Th{T}$ dan
!!a{derivation} atas~$\lnot !A$. Karena $\Th{T}$ lengkap, salah satunya pasti
ditemukan; dan karena $\Th{T}$ konsisten, jika !!a{derivation} atas~$\lnot !A$
ditemukan, tidak ada !!{derivation} atas~$!A$.

Dengan ungkapan lain, kita telah mengetahui bahwa $\Th{T}$ bersifat
!!{c.e.}; jadi, menurut suatu teorema yang telah dibuktikan sebelumnya, cukup
ditunjukkan bahwa komplemen~$\Th{T}$ juga bersifat !!{c.e.}. Namun,
!!a{formula}~$!A$ berada dalam~$\Th{\bar T}$ jika dan hanya jika $\lnot !A$
berada dalam~$\Th{T}$; jadi $\Th{\bar T} \leq_m \Th{T}$.
\end{proof}

\end{document}
